\documentclass[11pt]{article}

\usepackage[margin=1in]{geometry}
\usepackage{amsmath, amssymb, amsthm}
\usepackage{mathtools}
\usepackage{bm}
\usepackage{hyperref}
\usepackage{xcolor}
\usepackage{enumitem}

\newtheorem{theorem}{Theorem}
\newtheorem{lemma}[theorem]{Lemma}
\newtheorem{proposition}[theorem]{Proposition}
\newtheorem{corollary}[theorem]{Corollary}
\theoremstyle{definition}
\newtheorem{definition}[theorem]{Definition}
\newtheorem{assumption}[theorem]{Assumption}
\theoremstyle{remark}
\newtheorem{remark}[theorem]{Remark}

\newcommand{\R}{\mathbb{R}}
\newcommand{\E}{\mathbb{E}}
\newcommand{\Var}{\mathrm{Var}}
\newcommand{\Cov}{\mathrm{Cov}}
\newcommand{\KL}{\mathrm{KL}}
\newcommand{\TV}{\mathrm{TV}}
\newcommand{\D}{\mathcal{D}}
\newcommand{\Loss}{\mathcal{L}}
\newcommand{\eps}{\varepsilon}
\newcommand{\old}{\mathrm{old}}
\DeclareMathOperator*{\argmax}{arg\,max}

\title{Dropout-Induced Stochasticity for GRPO Fine-Tuning of \\
       Latent-Reasoning Language Models}
\author{Wooil Jung}
\date{\today}

\begin{document}
\maketitle

\begin{abstract}
Latent-reasoning language models such as \textsc{Coconut} (Hao et al., 2024)
replace discrete chain-of-thought tokens with continuous hidden states fed
recurrently into the next reasoning step.  Because the latent recurrence is
deterministic given the input, na\"ive policy-gradient methods such as GRPO
cannot generate diverse rollouts in latent space, which removes the variance
that group-relative advantage estimation relies on.  We propose to use a small
dropout rate $p \approx 0.1$ as the source of stochasticity during rollout, and
to \emph{replay the same dropout mask} during the policy update so that the
latent trajectory observed at rollout time is reproduced exactly.  We show
that, under this construction, dropout noise plays the role of an
\emph{exogenous} reparameterization variable, the per-mask GRPO surrogate is an
unbiased estimator of the policy gradient of an augmented (mask-marginalized)
policy, and replaying the mask provides common-random-numbers variance
reduction over the alternative of resampling at update time.  We also give an
inductive argument that the surrogate gradient is well defined at every latent
depth $T$, and we relate the construction to the MC-dropout Bayesian
interpretation of dropout.
\end{abstract}

% =====================================================================
\section{Introduction and motivation}
% =====================================================================

Group Relative Policy Optimization (GRPO; Shao et al., 2024) has become the
standard reinforcement-learning objective for fine-tuning reasoning LLMs on
math benchmarks.  Its appeal is that the per-prompt advantage is computed
purely from a group of $K$ rollouts and a verifier reward, with no learned
critic.  GRPO requires that the $K$ rollouts \emph{differ}, because the
group-normalized advantage,
\begin{equation}
A^{(k)} \;=\; \frac{r^{(k)} - \mu_r}{\sigma_r + \delta},
\qquad
\mu_r = \tfrac1K\sum_{k=1}^{K} r^{(k)},
\qquad
\sigma_r^2 = \tfrac1K\sum_{k=1}^{K}(r^{(k)} - \mu_r)^2,
\label{eq:grpo-advantage}
\end{equation}
is identically zero whenever all rollouts agree.  In a standard token-level
LLM, diversity is supplied automatically by sampling from the next-token
distribution.

\textsc{Coconut}-style latent reasoning, by contrast, does not sample.  Hidden
states $h_t \in \R^d$ are fed back into the model as the embedding of the next
``token'' for $t = 1, \dots, T$ until a stopping criterion fires; only after
the latent phase ends does the model emit discrete answer tokens
$y_1, \dots, y_L$.  Conditional on the prompt $x$, the latent phase is a
\emph{deterministic} function of the parameters $\theta$.  This breaks GRPO:
all $K$ rollouts share the same $h_T$ and (for greedy decoding) the same $y$.

We propose to inject stochasticity by enabling dropout during rollout.  This
yields per-rollout latent trajectories that depend on a Bernoulli mask
$\xi \sim p(\xi)$ which is independent of $\theta$, so GRPO can grade rollouts
relative to one another.  At update time we re-mount the \emph{same} mask so
that the gradient is computed against the trajectory that was actually graded.
The rest of this note formalizes the construction and proves the basic
correctness and variance properties.

% =====================================================================
\section{Preliminaries}
% =====================================================================

\subsection{Coconut recurrence with dropout}
\label{sec:coconut-recurrence}

Let $f_\theta : \R^d \to \R^d$ denote one Transformer step of the
latent-reasoning model with parameters $\theta \in \R^D$.  Standard
\textsc{Coconut} unrolls
\begin{equation}
h_t \;=\; f_\theta(h_{t-1}),
\qquad
h_0 \;=\; \mathrm{embed}_\theta(x),
\qquad t = 1,\dots,T.
\end{equation}
With dropout at rate $p$ and keep-probability $q := 1-p$, define a Bernoulli
mask $m_t \in \{0,1\}^D$ with $m_t \sim \mathrm{Bern}(q)^{\otimes D}$ and the
rescaled per-step parameters
\begin{equation}
\tilde\theta_t(\xi_t) \;:=\; \theta \odot m_t / q,
\qquad
\E_{\xi_t}\!\big[m_t/q\big] = \mathbf 1.
\label{eq:dropout-rescale}
\end{equation}
The dropout-perturbed recurrence is
\begin{equation}
h_t \;=\; f_{\tilde\theta_t(\xi_t)}(h_{t-1}),
\qquad
\xi := (\xi_1, \dots, \xi_T) \sim p(\xi) = \prod_{t=1}^T p(\xi_t).
\label{eq:dropout-recurrence}
\end{equation}
Once the latent phase finishes, answer tokens are sampled autoregressively,
\begin{equation}
y_\ell \;\sim\; \pi_\theta(\cdot \mid x, h_T, y_{<\ell}),
\qquad \ell = 1, \dots, L,
\end{equation}
optionally with dropout enabled here as well; we absorb any such sampling
randomness into $\xi$ by reparameterizing the categorical samples with
Gumbel noise when convenient.  A scalar reward $r(y, x) \in [0,1]$ is provided
by an external verifier (e.g.\ exact-match on \textsc{GSM8K}).

\subsection{Augmented policy}
The induced distribution over $(\xi, y)$ given $x$ is
\begin{equation}
\Pi_\theta(\xi, y \mid x)
   \;=\; p(\xi) \, \pi_\theta(y \mid x, \xi),
\qquad
\pi_\theta(y \mid x, \xi)
  := \prod_{\ell=1}^{L} \pi_\theta\!\big(y_\ell \,\big|\, x, h_T(\theta,\xi,x), y_{<\ell}\big).
\label{eq:augmented-policy}
\end{equation}
We treat $\Pi_\theta$ as the policy that GRPO optimizes; the dropout mask is
part of the action.  The marginal policy
\begin{equation}
\bar\pi_\theta(y \mid x) \;:=\; \E_{\xi \sim p}\!\big[\pi_\theta(y \mid x, \xi)\big]
\label{eq:marginal-policy}
\end{equation}
is the object whose expected reward we ultimately wish to maximize.

\subsection{GRPO objective}

For a group of $K$ rollouts $\{(\xi^{(k)}, y^{(k)})\}_{k=1}^K$ generated under
parameters $\theta_\old$, with rewards $r^{(k)} = r(y^{(k)}, x)$, define
advantages $A^{(k)}$ via~\eqref{eq:grpo-advantage} and the per-token
importance ratio
\begin{equation}
\rho_\theta^{(k)}(\xi, y)
  \;:=\;
\frac{\pi_\theta(y \mid x, \xi)}{\pi_{\theta_\old}(y \mid x, \xi)}.
\end{equation}
The clipped GRPO surrogate is
\begin{equation}
\Loss^{\mathrm{GRPO}}(\theta)
  \;=\;
\E_{x \sim \D}\,
\frac{1}{K}\sum_{k=1}^{K}
  \min\!\Big(
    \rho_\theta^{(k)}\, A^{(k)}, \;
    \mathrm{clip}(\rho_\theta^{(k)},\,1-\eps,\,1+\eps)\,A^{(k)}
  \Big)
  \;-\;\beta\,\KL\!\big(\pi_{\theta}\,\big\|\,\pi_{\mathrm{ref}}\big).
\label{eq:grpo-loss}
\end{equation}

% =====================================================================
\section{Method: dropout as exogenous noise with mask replay}
% =====================================================================

The proposed procedure is:
\begin{enumerate}[topsep=2pt,itemsep=2pt]
\item For each prompt $x$, draw $K$ independent dropout masks
      $\xi^{(1)}, \dots, \xi^{(K)} \sim p(\xi)$.
\item For each $k$, run the latent-reasoning recurrence
      \eqref{eq:dropout-recurrence} with mask $\xi^{(k)}$ \emph{stored on
      device}, decode answer tokens $y^{(k)}$, compute reward $r^{(k)}$ and
      advantage $A^{(k)}$ via~\eqref{eq:grpo-advantage}.
\item During the policy update, recompute $\pi_\theta(y^{(k)}\mid x,\xi^{(k)})$
      with the \emph{same} stored mask $\xi^{(k)}$, so that $h_T$ at update
      time matches $h_T$ at rollout time exactly when $\theta = \theta_\old$.
      Form the GRPO loss~\eqref{eq:grpo-loss} and step $\theta$.
\end{enumerate}

The construction rests on three observations.

\paragraph{(i) Independence of $\xi$ from $\theta$.}
The dropout mask $\xi$ is sampled from a fixed distribution $p(\xi)$ that does
not depend on $\theta$.  Hence, in~\eqref{eq:augmented-policy},
$\nabla_\theta \log p(\xi) = 0$, and
\begin{equation}
\nabla_\theta \log \Pi_\theta(\xi, y \mid x)
  \;=\;
\nabla_\theta \log \pi_\theta(y \mid x, \xi).
\label{eq:score-reduces}
\end{equation}
The gradient signal therefore comes entirely from the $\xi$-conditional log
likelihood of the answer tokens, which is differentiable through the
unrolled latent path.

\paragraph{(ii) Reparameterization view.}
With $\xi$ fixed, $h_T = h_T(\theta, \xi, x)$ is a deterministic, almost
everywhere differentiable function of $\theta$.  This is precisely the setting
of the reparameterization trick (Kingma \& Welling, 2014).  Equivalently,
sampling $\xi$ first and then computing $h_T$ is a coupling of stochastic
processes across different $\theta$ values --- the common-random-numbers
construction (Glasserman \& Yao, 1992).

\paragraph{(iii) Mask replay = common random numbers.}
At update time we evaluate $\pi_\theta$ on the same $\xi^{(k)}$ that was used
at rollout time.  This couples the rollout-time and update-time trajectories
so that the only source of difference between them is the parameter update
$\theta - \theta_\old$ itself.  Without replay, the importance ratio
$\rho_\theta(\xi, y) / \rho_\theta(\xi', y)$ would also reflect mask
resampling, inflating its variance.

% =====================================================================
\section{Theoretical analysis}
% =====================================================================

\subsection{Well-definedness of the surrogate gradient (induction on $T$)}

\begin{assumption}\label{ass:smooth}
For every fixed mask $\xi$, the map $\theta \mapsto f_{\tilde\theta(\xi)}(h)$
is differentiable in $\theta$ for almost every $h$, and the answer-token
log-likelihood $\log \pi_\theta(y \mid x, h_T, y_{<\ell})$ is differentiable
in both $\theta$ and $h_T$.
\end{assumption}

This is satisfied by every standard Transformer block: the dropout rescaling
$\theta \odot m/q$ is linear in $\theta$, and the remaining nonlinearities
(LayerNorm, GELU, softmax) are almost everywhere differentiable.

\begin{lemma}[Well-defined latent gradient at depth $T$]
\label{lem:induction}
Under Assumption~\ref{ass:smooth}, for every $T \ge 0$ and almost every
$\xi$, the Jacobian $\partial h_T / \partial \theta$ exists and is given by
the recursion
\begin{equation}
\frac{\partial h_t}{\partial \theta}
  \;=\;
\frac{\partial f_{\tilde\theta_t}}{\partial \theta}\bigg|_{h_{t-1}}
  \cdot
\frac{\partial \tilde\theta_t}{\partial \theta}
  \;+\;
\frac{\partial f_{\tilde\theta_t}}{\partial h}\bigg|_{h_{t-1}}
  \cdot
\frac{\partial h_{t-1}}{\partial \theta},
\qquad
\frac{\partial h_0}{\partial \theta}
  \;=\;
\frac{\partial \,\mathrm{embed}_\theta(x)}{\partial \theta},
\label{eq:latent-jacobian}
\end{equation}
where $\partial \tilde\theta_t / \partial \theta = \mathrm{diag}(m_t)/q$ is the
dropout-rescaled identity.  Consequently
$\nabla_\theta \log \pi_\theta(y \mid x, \xi)$ is well defined.
\end{lemma}

\begin{proof}
By induction on $t$.

\emph{Base case} ($t=0$).  $h_0 = \mathrm{embed}_\theta(x)$ is differentiable in
$\theta$ by assumption, giving the stated initial Jacobian.

\emph{Inductive step.}  Assume $\partial h_{t-1}/\partial \theta$ exists.
Write $h_t = f(\tilde\theta_t, h_{t-1})$, viewing $f$ as a function of two
arguments.  By the multivariate chain rule and Assumption~\ref{ass:smooth},
\[
\frac{\partial h_t}{\partial \theta}
  \;=\;
\frac{\partial f}{\partial \tilde\theta_t}
  \cdot
\frac{\partial \tilde\theta_t}{\partial \theta}
  \;+\;
\frac{\partial f}{\partial h_{t-1}}
  \cdot
\frac{\partial h_{t-1}}{\partial \theta},
\]
and $\partial \tilde\theta_t/\partial \theta = \mathrm{diag}(m_t)/q$ from
\eqref{eq:dropout-rescale}.  Both terms exist (the second by the inductive
hypothesis), giving~\eqref{eq:latent-jacobian}.

Finally, since $\log \pi_\theta(y \mid x, h_T, y_{<\ell})$ is differentiable in
both $\theta$ and $h_T$, one more chain-rule application yields
$\nabla_\theta \log \pi_\theta(y \mid x, \xi)
   = \sum_\ell [\partial_\theta \log \pi_\theta]_{\text{direct}}
              + [\partial_{h_T}\log \pi_\theta]\cdot \partial h_T/\partial\theta$.
\end{proof}

\begin{remark}
The recursion~\eqref{eq:latent-jacobian} is what \texttt{loss.backward()}
implements when the dropout masks are stored and replayed: PyTorch's
\texttt{torch.nn.functional.dropout} treats the rescaled mask as a constant
during backprop, which is exactly $\partial \tilde\theta_t/\partial \theta =
\mathrm{diag}(m_t)/q$.
\end{remark}

\subsection{Unbiasedness of the dropout-GRPO surrogate}

\begin{theorem}[Unbiased policy gradient on the augmented policy]
\label{thm:unbiased}
Let $J(\theta) := \E_{x\sim\D}\,\E_{(\xi,y)\sim\Pi_\theta(\cdot\mid x)}[r(y,x)]$
be the expected reward under the augmented policy~\eqref{eq:augmented-policy}.
Then
\begin{equation}
\nabla_\theta \Loss^{\mathrm{GRPO}}(\theta)\Big|_{\theta=\theta_\old}
  \;=\;
\E_{x,\xi,y}\!\left[A(\xi, y)\,
       \nabla_\theta \log \pi_\theta(y\mid x,\xi)\right]_{\theta_\old}
  \;=\;
\nabla_\theta J(\theta)\Big|_{\theta_\old},
\label{eq:unbiased}
\end{equation}
provided the KL term in~\eqref{eq:grpo-loss} is evaluated as
$\KL(\Pi_\theta \| \Pi_\mathrm{ref})$ on the augmented policy.
\end{theorem}

\begin{proof}
At $\theta = \theta_\old$ we have $\rho_{\theta_\old} \equiv 1$, so the clip is
inactive and the gradient of the $\min$ reduces to the gradient of the unclipped
ratio:
\[
\nabla_\theta \rho_\theta\big|_{\theta_\old}
   \;=\;
\nabla_\theta \log \pi_\theta(y\mid x,\xi)\big|_{\theta_\old}.
\]
Therefore
\begin{equation}
\nabla_\theta \Loss^{\mathrm{GRPO}}\big|_{\theta_\old}
  \;=\;
\E_{x,\xi,y}\!\left[A(\xi,y)\,\nabla_\theta\log\pi_\theta(y\mid x,\xi)\right]_{\theta_\old}.
\label{eq:proof-step1}
\end{equation}
By~\eqref{eq:score-reduces},
$\nabla_\theta\log\pi_\theta(y\mid x,\xi)
   = \nabla_\theta\log \Pi_\theta(\xi,y\mid x)$, so the right-hand side is the
REINFORCE estimator on the augmented policy with advantage $A$.  The
group-relative baseline $\mu_r$ is independent of any single rollout's score
function in expectation (a standard control-variate identity), so it
contributes zero bias.  Hence the right-hand side equals
$\nabla_\theta J(\theta)|_{\theta_\old}$.
\end{proof}

\begin{corollary}[Bias on the marginal policy]
The estimator is unbiased for $J(\theta)$, i.e.\ for the expected reward of
the mask-marginalized policy~\eqref{eq:marginal-policy}, since
$\E_\xi[\pi_\theta(y\mid x,\xi)] = \bar\pi_\theta(y\mid x)$ implies
$\E_{x,\xi,y\sim\Pi_\theta}[r] = \E_{x,y\sim\bar\pi_\theta}[r]$.  In other
words, optimizing the augmented policy is exactly optimizing expected reward
under the dropout ensemble.
\end{corollary}

\subsection{Variance reduction via mask replay (CRN)}

Consider two estimators of $\nabla_\theta J$ at a perturbed parameter
$\theta = \theta_\old + \Delta$:
\begin{align}
\hat g_{\mathrm{CRN}} &:=
   A(\xi, y)\,\nabla_\theta \log \pi_\theta(y\mid x, \xi),
   \quad \xi\text{ shared with rollout,}
\\
\hat g_{\mathrm{indep}} &:=
   A(\xi, y)\,\nabla_\theta \log \pi_\theta(y\mid x, \xi'),
   \quad \xi' \perp\!\!\!\perp \xi\text{ resampled.}
\end{align}

\begin{proposition}[CRN variance reduction]
\label{prop:crn}
Suppose the trajectory-conditional reward $\bar r(\xi) := \E_y[r \mid \xi]$ and
the score function $s(\xi, y; \theta) := \nabla_\theta \log\pi_\theta(y\mid x,\xi)$
are positively correlated across $\xi$, i.e.\
$\Cov_\xi(\bar r(\xi),\,s(\xi,y;\theta)) \succeq 0$ component-wise.  Then
\begin{equation}
\Var[\hat g_{\mathrm{CRN}}] \;\le\; \Var[\hat g_{\mathrm{indep}}].
\end{equation}
\end{proposition}

\begin{proof}[Proof sketch]
Both estimators have the same mean (Theorem~\ref{thm:unbiased}).  Decompose
each variance into a within-mask and across-mask component using the law of
total variance.  The within-mask components are equal.  For the across-mask
component, the resampling estimator $\hat g_{\mathrm{indep}}$ has uncorrelated
$\bar r(\xi)$ and $s(\xi',y;\theta)$ (since $\xi'\perp\xi$), whereas
$\hat g_{\mathrm{CRN}}$ has covariance equal to the assumed nonnegative
quantity.  Subtracting yields the bound.  This is the standard CRN argument
(Glasserman \& Yao, 1992, Thm 3.1) specialized to the score-function
estimator.
\end{proof}

\begin{remark}
The positive-correlation assumption is mild in practice: masks that yield a
correct answer ($\bar r(\xi)$ large) also yield a score function that pushes
$\theta$ in directions that increase $\pi_\theta$ on those tokens.  Empirical
verification can be done by tracking the sign of the diagonal of
$\Cov_\xi(\bar r, s)$ during training.
\end{remark}

\subsection{Connection to MC-dropout}

Following Gal \& Ghahramani (2016), interpret each mask as a posterior sample
$\theta_\xi := \theta\odot m(\xi)/q$ from a variational distribution
$q_\theta(\theta')$ over parameters.  The marginal policy
\eqref{eq:marginal-policy} is then
\begin{equation}
\bar\pi_\theta(y\mid x)
  \;=\;
\E_{\theta'\sim q_\theta}[\pi_{\theta'}(y\mid x)],
\end{equation}
i.e.\ a Bayesian-model-average policy.  Theorem~\ref{thm:unbiased} states
that dropout-GRPO performs ensemble-aware policy optimization: each gradient
step is the Monte Carlo average of per-ensemble-member policy gradients,
with the group-normalized advantage acting as a control variate.

% =====================================================================
\section{Algorithm}
% =====================================================================

\begin{figure}[h]
\centering
\fbox{\begin{minipage}{0.92\linewidth}
\textbf{Algorithm 1: Dropout-GRPO for latent-reasoning LLMs.}\\[2pt]
\textit{Inputs:} base policy $\pi_{\theta_\old}$, prompts $\D$, group size
$K$, dropout rate $p$, latent depth $T$, clip $\eps$, KL coef $\beta$,
mini-epochs $E$.
\begin{enumerate}[topsep=2pt,itemsep=1pt,leftmargin=*]
\item For each prompt $x \sim \D$:
  \begin{enumerate}[topsep=1pt,itemsep=1pt,label=\alph*.]
  \item \emph{Rollout} (train mode, dropout on). For $k=1,\dots,K$:
    sample fresh dropout RNG seeds, store seeds for $\xi^{(k)}$, unroll
    latent recurrence \eqref{eq:dropout-recurrence}, decode $y^{(k)}$,
    record per-token logprobs $\log\pi_{\theta_\old}(y^{(k)}\mid x,\xi^{(k)})$,
    compute $r^{(k)} = \mathrm{verifier}(y^{(k)}, x)$.
  \item Compute advantages
        $A^{(k)} = (r^{(k)} - \mu_r)/(\sigma_r + \delta)$
        as in \eqref{eq:grpo-advantage}.
  \item \emph{Update} (mask replay). For mini-epoch $e=1,\dots,E$ and
        each $k$: regenerate $\xi^{(k)}$ from stored seeds, recompute
        $\log\pi_\theta(y^{(k)}\mid x, \xi^{(k)})$, form
        $\rho^{(k)} = \pi_\theta/\pi_{\theta_\old}$, take a gradient step
        on $\Loss^{\mathrm{GRPO}}$ from \eqref{eq:grpo-loss}.
  \item Set $\theta_\old \gets \theta$.
  \end{enumerate}
\end{enumerate}
\end{minipage}}
\label{alg:dropout-grpo}
\end{figure}

% =====================================================================
\section{Practical concerns and open questions}
% =====================================================================

\paragraph{Sensitivity of $h_T$ to $\xi$.}
If dropout perturbations to $h_T$ do not change the argmax of the answer-token
distribution, all $K$ rollouts agree, advantages collapse to zero, and no
gradient is produced.  This is the dropout-rate analogue of temperature in
token-level RL.  A diagnostic to monitor is the
within-group standard deviation of $r^{(k)}$ (Eq.~\eqref{eq:grpo-advantage});
if $\sigma_r \to 0$, raise $p$ or fold token-sampling Gumbel noise into $\xi$.

\paragraph{Importance-ratio explosion.}
The ratio $\rho_\theta(\xi,y)$ couples parameter updates and dropout-induced
logit shifts, and can be large because dropping different units changes logits
non-smoothly.  The PPO-style clip $\eps$ may need to be tightened relative to
the token-LM regime, or applied per latent step rather than only per answer
token.

\paragraph{Validity of clip under CRN.}
The clip operator in~\eqref{eq:grpo-loss} is non-smooth, so
Theorem~\ref{thm:unbiased} as stated holds only at $\theta = \theta_\old$
(initial step).  Standard PPO analysis (Schulman et al., 2017, \S 5) extends
the bound to small steps under a Lipschitz condition on $\rho_\theta$, which
in turn requires $\rho_\theta$ to remain bounded under mask replay; this is
plausible but should be verified empirically (track max ratio).

\paragraph{Train/deploy mismatch.}
Inference typically disables dropout, so the deployed policy is
$\pi_\theta(y\mid x, \xi=\mathbf 1)$ rather than the marginal
$\bar\pi_\theta$.  Two options: (a)~keep dropout on at deploy time and average
over masks, or (b)~accept a bias bounded by
$\TV(\bar\pi_\theta,\,\pi_\theta(\cdot\mid \mathbf 1)) = \mathcal{O}(p\sqrt d)$
under standard regularity assumptions on the network's Lipschitz constant in
its mask argument.

\paragraph{Memory for stored masks.}
Storing $\xi^{(k)}$ for $K$ rollouts and $T$ latent steps over a $D$-parameter
model costs $\mathcal{O}(KTD)$ bits.  In practice, only the dropout layers'
masks need to be stored (not full-parameter masks), and a per-layer RNG seed
is sufficient: store the seed, regenerate the mask deterministically at
update time.  This is the recommended implementation.

% =====================================================================
\section{Summary}
% =====================================================================

\begin{itemize}[topsep=2pt,itemsep=2pt]
\item Latent-reasoning models break GRPO because the latent phase is
      deterministic; dropout supplies the missing stochasticity without
      changing the model class.
\item Dropout noise is exogenous (independent of $\theta$), so it acts as a
      reparameterization variable; the score function reduces to the
      $\xi$-conditional log likelihood (Eq.~\ref{eq:score-reduces}).
\item The latent gradient $\partial h_T/\partial \theta$ is well defined at
      every depth $T$ (Lemma~\ref{lem:induction}), justifying backprop through
      the unrolled latent recurrence.
\item GRPO with replayed masks is an unbiased estimator of the policy
      gradient on the dropout-augmented policy, hence on the marginal
      policy (Theorem~\ref{thm:unbiased}).
\item Replaying the rollout-time mask at update time gives a CRN coupling
      that lower-bounds variance against the resample-at-update alternative
      (Proposition~\ref{prop:crn}).
\item Practical caveats are: $h_T$-sensitivity, importance-ratio control,
      clip-region Lipschitz behavior, and train/deploy mismatch.
\end{itemize}

% =====================================================================
\section*{References (informal)}
% =====================================================================

\begin{itemize}[topsep=2pt,itemsep=1pt,leftmargin=*]
\item Hao, S.\ et al.\ (2024). \emph{Training large language models to reason
      in a continuous latent space.}  (Coconut.)
\item Shao, Z.\ et al.\ (2024).  \emph{DeepSeekMath: Pushing the limits of
      mathematical reasoning in open language models.} (GRPO.)
\item Schulman, J.\ et al.\ (2017).  \emph{Proximal policy optimization
      algorithms.}
\item Gal, Y.\ \& Ghahramani, Z.\ (2016).  \emph{Dropout as a Bayesian
      approximation: Representing model uncertainty in deep learning.}
\item Kingma, D.\ \& Welling, M.\ (2014).  \emph{Auto-encoding variational
      Bayes.}
\item Glasserman, P.\ \& Yao, D.D.\ (1992).  \emph{Some guidelines and
      guarantees for common random numbers.}
\end{itemize}

\end{document}
